\documentclass[11pt]{article}
\topmargin=-1.0cm
\usepackage{graphicx}
\usepackage{epstopdf}
\usepackage{amssymb}
\usepackage{amsmath}
\usepackage{amsthm}
\usepackage{color}
\usepackage[margin=1in]{geometry}
\usepackage[hidelinks]{hyperref}
\usepackage{booktabs}

% ******************************* Author Macros *******************************

\newcommand{\pl}[2]{\frac{\partial#1}{\partial#2}}
\newcommand{\bu}{{\bf u} }
\newcommand{\p}{\partial}
\newcommand{\og}{\omega}
\newcommand{\Og}{\Omega}
\newcommand{\fl}[2]{\frac{#1}{#2}}
\newcommand{\dt}{\delta}
\newcommand{\nn}{\nonumber}
\newcommand{\ap}{\alpha}
\newcommand{\bt}{\beta}
\newcommand{\tht}{\theta}
\newcommand{\kp}{\kappa}
\newcommand{\Dt}{\Delta}
\renewcommand{\theequation}{\arabic{section}.\arabic{equation}}
\newcommand{\be}{\begin{equation}}
\newcommand{\ee}{\end{equation}}
\newcommand{\ba}{\begin{array}}
\newcommand{\ea}{\end{array}}
\newcommand{\bea}{\begin{eqnarray}}
\newcommand{\eea}{\end{eqnarray}}
\newcommand{\beas}{\begin{eqnarray*}}
\newcommand{\eeas}{\end{eqnarray*}}
\newtheorem{theorem}{Theorem}[section]
\newtheorem{lemma}{Lemma}[section]
\newtheorem{proposition}{Proposition}[section]
\newtheorem{remark}{Remark}[section]
\newtheorem{assumption}{Assumption}[section]
\newcommand{\bx}{{\bf x} }
\newcommand{\gm}{\gamma}
\newcommand{\vep}{\varepsilon}
\newcommand{\T}{\mathbb T}
\newcommand{\dT}{d_{\mathbb T}}

% *****************************************************************************

\title{Appendix deleted}
\begin{document}
\date{}
\maketitle

\section{Numerical Hamiltonians and one-sided derivative reconstruction}\label{app:hamiltonian_weno}

For the monotone theory in Section~3, the numerical Hamiltonian \(\widehat H(a,b)\) is assumed to be consistent with \(-p^2\), nondecreasing in its first argument, and nonincreasing in its second argument.  Two standard choices are the Lax--Friedrichs Hamiltonian
\begin{equation}\label{eq:app_LF_H_new}
    \widehat H_{\rm LF}(a,b)
    =-\left(\frac{a+b}{2}\right)^2+\frac{\alpha}{2}(a-b),
\end{equation}
where \(\alpha\) is chosen no smaller than the relevant characteristic speed, and the Godunov Hamiltonian
\begin{equation}\label{eq:app_Godunov_H_new}
    \widehat H_G(a,b)=-(a^-)^2-(b^+)^2,
    \qquad
    a^-:=\min\{a,0\},
    \qquad
    b^+:=\max\{b,0\}.
\end{equation}
The stationary phase estimate in Proposition~\ref{prop:stationary_phase_estimate} uses \(\widehat H_G\).

Some numerical runs use a high-order one-sided WENO reconstruction for the derivatives entering the Hamiltonian.  This is an implementation option, not the monotone Hamiltonian used in the proof.  For a periodic grid function \(u_k\), define the left-biased quantities
\[
    a=u_{k-3}-2u_{k-2}+u_{k-1},\quad
    b=u_{k-2}-2u_{k-1}+u_k,
\]
\[
    c=u_{k-1}-2u_k+u_{k+1},\quad
    d=u_k-2u_{k+1}+u_{k+2}.
\]
Set
\[
    I_0=13(a-b)^2+3(a-3b)^2,
    \quad
    I_1=13(b-c)^2+3(b+c)^2,
\]
\[
    I_2=13(c-d)^2+3(3c-d)^2,
\]
and
\[
    \alpha_0=\frac{1}{(\epsilon_0+I_0)^2},
    \quad
    \alpha_1=\frac{6}{(\epsilon_0+I_1)^2},
    \quad
    \alpha_2=\frac{3}{(\epsilon_0+I_2)^2},
    \quad
    \omega_r=\frac{\alpha_r}{\alpha_0+\alpha_1+\alpha_2}.
\]
The left derivative reconstruction used in the implementation is
\begin{equation}\label{eq:weno_left_derivative_new}
\begin{aligned}
    p_k^-=\frac{1}{\Delta\theta}
    \bigg[&\frac{-(u_{k-1}-u_{k-2})+7(u_k-u_{k-1})+7(u_{k+1}-u_k)-(u_{k+2}-u_{k+1})}{12} \\
    &-\frac{\omega_0}{3}(a-2b+c)
    -\frac{\omega_2-1/2}{6}(b-2c+d)
    \bigg].
\end{aligned}
\end{equation}
The right-biased derivative \(p_k^+\) is obtained by applying the same formula to the reversed stencil.  Equivalently, with
\[
    a^+=u_{k+3}-2u_{k+2}+u_{k+1},\quad
    b^+=u_{k+2}-2u_{k+1}+u_k,
\]
\[
    c^+=u_{k+1}-2u_k+u_{k-1},\quad
    d^+=u_k-2u_{k-1}+u_{k-2},
\]
and the corresponding weights \(\omega_0^+,\omega_2^+\),
\begin{equation}\label{eq:weno_right_derivative_new}
\begin{aligned}
    p_k^+=\frac{1}{\Delta\theta}
    \bigg[&\frac{-(u_{k-1}-u_{k-2})+7(u_k-u_{k-1})+7(u_{k+1}-u_k)-(u_{k+2}-u_{k+1})}{12} \\
    &+\frac{\omega_0^+}{3}(a^+-2b^++c^+)
    +\frac{\omega_2^+-1/2}{6}(b^+-2c^++d^+)
    \bigg].
\end{aligned}
\end{equation}
The WENO option then uses a Roe-type entropy-fix numerical Hamiltonian for \(H(p)=-p^2\) with the reconstructed values \((p_k^+,p_k^-)\).  Since this high-order reconstruction is not monotone in the sense assumed in Section~3, the numerical experiments report it separately from the monotone Godunov/Lax--Friedrichs choices when needed.

\section{Conservative fluxes for the split amplitude equation}\label{app:fluxes_new}

The continuous conservative flux in the split amplitude equation is
\[
    F=-W\partial_\theta u.
\]
On a uniform periodic trait grid, write
\[
    D_\theta\widehat F(W,u)_{j,k}
    =\frac{\widehat F_{j,k+1/2}-\widehat F_{j,k-1/2}}{\Delta\theta},
    \qquad
    a_{k+1/2}=-\frac{u_{k+1}-u_k}{\Delta\theta}.
\]
The upwind flux is
\begin{equation}\label{eq:app_upwind_flux_new}
    \widehat F^{\rm up}_{j,k+1/2}
    =a_{k+1/2}^+ W_{j,k}+a_{k+1/2}^- W_{j,k+1},
    \qquad
    a^+=\max\{a,0\},
    \qquad
    a^-=\min\{a,0\}.
\end{equation}
A Lax--Friedrichs flux is
\begin{equation}\label{eq:app_LF_flux_new}
    \widehat F^{\rm LF}_{j,k+1/2}
    =\frac{1}{2}a_{k+1/2}(W_{j,k}+W_{j,k+1})
    -\frac{1}{2}\alpha_{k+1/2}(W_{j,k+1}-W_{j,k}),
    \qquad
    \alpha_{k+1/2}\ge |a_{k+1/2}|.
\end{equation}
Both fluxes are conservative in the periodic trait variable.  The upwind flux also satisfies the quasi-positivity condition used in Lemma~\ref{lem:positivity_W}.






\section{Scalar mass-balance correction}
\label{app:mass_correction}

A detailed algebraic form of the scalar mass-balance correction used in the
fully discrete method is provided here.  
All spatial integrals below are understood in the discrete sense and may be
evaluated by any consistent spatial quadrature used in the fully discrete
scheme.
The correction is based on the mass identity
\[
    \varepsilon \frac{d}{dt}\int_\Omega \rho(x,t)\,dx
    =
    \int_\Omega \rho(x,t)\bigl(K(x)-\rho(x,t)\bigr)\,dx .
\]
Let \(\widetilde W^{n+1}\), \(\widetilde u^{n+1}\), and
\(\widetilde\rho^{n+1}\) denote the provisional amplitude, phase, and
reconstructed spatial marginal after one time step.  The correction is defined
by
\[
    W^{n+1}=\lambda \widetilde W^{n+1},
    \qquad
    u^{n+1}=\widetilde u^{n+1},
    \qquad
    \lambda>0 .
\]
It follows that
\[
    \rho^{n+1}=\lambda\widetilde\rho^{n+1},
    \qquad
    n^{n+1}=\lambda\widetilde n^{n+1},
    \qquad
    P^{n+1}=\lambda\widetilde P^{n+1}.
\]
Thus the correction changes only the density scale and leaves the phase profile
unchanged.


For compactness, define
\[
    \widetilde M
    =
    \int_\Omega \widetilde\rho^{n+1}\,dx,
    \qquad
    \widetilde S
    =
    \int_\Omega K\widetilde\rho^{n+1}\,dx,
    \qquad
    \widetilde Q
    =
    \int_\Omega (\widetilde\rho^{n+1})^2\,dx,
\]
and
\[
    M^n
    =
    \int_\Omega \rho^n\,dx,
    \qquad
    R^n
    =
    \int_\Omega \rho^n(K-\rho^n)\,dx .
\]
If the trapezoidal discretization is applied in time, the correction factor
\(\lambda\) is determined by
\[
    \varepsilon
    \frac{\lambda\widetilde M-M^n}{\tau}
    =
    \frac12
    \left[
        R^n+\lambda\widetilde S-\lambda^2\widetilde Q
    \right].
\]
Equivalently, \(\lambda\) is the positive root of
\[
    \mathcal A_{\rm TR}\lambda^2
    +
    \mathcal B_{\rm TR}\lambda
    +
    \mathcal C_{\rm TR}
    =
    0,
\]
where
\[
    \mathcal A_{\rm TR}
    =
    \frac{\tau}{2}\widetilde Q,
    \qquad
    \mathcal B_{\rm TR}
    =
    \varepsilon\widetilde M
    -
    \frac{\tau}{2}\widetilde S,
    \qquad
    \mathcal C_{\rm TR}
    =
    -\varepsilon M^n
    -
    \frac{\tau}{2}R^n .
\]
If a backward Euler discretization of
the reaction term is applied,  \(\lambda\) is determined by
\[
    \varepsilon
    \frac{\lambda\widetilde M-M^n}{\tau}
    =
    \lambda\widetilde S-\lambda^2\widetilde Q,
\]
or equivalently
\[
    \mathcal A_{\rm BE}\lambda^2
    +
    \mathcal B_{\rm BE}\lambda
    +
    \mathcal C_{\rm BE}
    =
    0,
\]
with
\[
    \mathcal A_{\rm BE}
    =
    \tau\widetilde Q,
    \qquad
    \mathcal B_{\rm BE}
    =
    \varepsilon\widetilde M
    -
    \tau\widetilde S,
    \qquad
    \mathcal C_{\rm BE}
    =
    -\varepsilon M^n .
\]
The positive root is selected in both cases to preserve positivity of the
reconstructed density.


When the provisional reconstruction already satisfies the discrete mass balance,
the correction factor is close to one.  In the small mutation regime, however,
a small mass inconsistency may lead to a visible error in the peak height of the
trait marginal.  The scalar correction therefore acts as a mass-balance
projection for the reconstructed density-level quantities.  It improves the
stability of the total mass and the peak height without changing the
phase-based peak location.



\section{Density and trait-marginal reconstruction, gauge normalization, and dual trait grids}\label{app:reconstruction}

The reconstruction operators used in the code act on the WKB variables and do
not solve the original density equation.  The default safe reconstruction of
\(\rho\) is the log-stabilized quadrature.  With
\[
    \ell_{j,k}=\log\bigl(\max\{W_{j,k},W_{\rm floor}\}\bigr)
    +\frac{u_k}{\varepsilon},
\]
the code evaluates
\[
    \rho_j=\sum_k\omega_k\exp(\ell_{j,k})
\]
by a row-wise log-sum-exp formula.  This is algebraically identical to the
composite quadrature
\[
    \rho_{j,Q}=\Delta\theta\sum_k W_{j,k}\exp(u_k/\varepsilon),
\]
but it avoids underflow and overflow when \(\varepsilon\) is small.  The
max-gauge \eqref{eq:gauge_new} is applied after each WKB step.  It leaves
\(n_{j,k}=W_{j,k}\exp(u_k/\varepsilon)\) unchanged and keeps
\(\max_k u_k=0\).

A Laplace reconstruction can be used near a nondegenerate phase maximizer.  Let
\(k_*\) be a nodal maximizer of \(u_k\).  The periodic stencil is represented in
local coordinates
\[
    z_{k_*-1}=-\Delta\theta,
    \qquad
    z_{k_*}=0,
    \qquad
    z_{k_*+1}=\Delta\theta,
\]
even when the stencil crosses the endpoint of \([0,1)\).  With
\[
    d_{\theta\theta}u_{k_*}
    =\frac{u_{k_*+1}-2u_{k_*}+u_{k_*-1}}{(\Delta\theta)^2},
\]
the parabolic vertex is
\begin{equation}\label{eq:laplace_vertex_new}
    z_*=-\frac{\Delta\theta}{2}
    \frac{u_{k_*+1}-u_{k_*-1}}{u_{k_*+1}-2u_{k_*}+u_{k_*-1}}.
\end{equation}
The reconstructed trait location is
\(\theta_*=(\theta_{k_*}+z_*)\bmod 1\).  Interpolating \(W\) and \(u\)
quadratically on this local stencil gives \(W(x_j,\theta_*)\) and
\(u(\theta_*)\), and
\begin{equation}\label{eq:laplace_rho_new}
    \rho_j^{\rm Lap}
    =W(x_j,\theta_*)\exp\left(\frac{u(\theta_*)}{\varepsilon}\right)
    \left(\frac{2\pi\varepsilon}{|d_{\theta\theta}u_{k_*}|}\right)^{1/2}.
\end{equation}
In the implementation this expression is evaluated in logarithmic form.  If
\(d_{\theta\theta}u_{k_*}\ge0\), the curvature is too small, or the interpolated
amplitude is non-positive, the hybrid reconstruction falls back to the
log-stabilized direct quadrature.

The trait marginal is reconstructed differently from \(\rho\).  For WKB data,
the code first forms the amplitude marginal
\[
    A_k=\int_\Omega W(x,\theta_k)\,dx
\]
and then works with
\begin{equation}\label{eq:appendix_logP_new}
    L_k=\log P_k=\frac{u_k}{\varepsilon}+\log\bigl(\max\{A_k,A_{\rm floor}\}\bigr).
\end{equation}
Thus interpolation is applied to \(L=\log P\), not to \(P\) itself.  The peak
search uses a periodic WENO interpolation of \(L_k\) on a fine uniform search
grid, followed by the three-point parabolic refinement \eqref{eq:laplace_vertex_new}
applied to the reconstructed values on that search grid.  This is why the
reported small-\(\varepsilon\) peak location is less sensitive to whether the
true density peak lies between two coarse trait nodes.

After the peak location \(\theta_*\) has been found, the code can reconstruct a
local peak shape.  In the phase-amplitude form, it fits
\[
    u(\theta_*+z)\simeq \alpha_0+\alpha_1z+\alpha_2z^2,
    \qquad
    \log A(\theta_*+z)\simeq \beta_0+\beta_1z+\beta_2z^2
\]
on a small periodic stencil, optionally with weights that emphasize points near
the peak.  Hence
\[
    \log P(\theta_*+z)
    \simeq
    \left(\frac{\alpha_0}{\varepsilon}+\beta_0\right)
    +\left(\frac{\alpha_1}{\varepsilon}+\beta_1\right)z
    +\left(\frac{\alpha_2}{\varepsilon}+\beta_2\right)z^2 .
\]
When the quadratic coefficient is negative, the fitted peak width is
\begin{equation}\label{eq:appendix_peak_width_new}
    \sigma_P\simeq \frac{1}{\sqrt{-2(\alpha_2/\varepsilon+\beta_2)}}.
\end{equation}
The reconstructed shape is then normalized by a prescribed mass,
\begin{equation}\label{eq:appendix_mass_shape_new}
    P_{\rm rec}(\theta)
    =M_{\rm target}
    \frac{\exp(\widehat L(\theta)-\max_\eta\widehat L(\eta))}
    {\int_{\mathbb T}\exp(\widehat L(\eta)-\max_\xi\widehat L(\xi))\,d\eta},
\end{equation}
where \(M_{\rm target}\) is the solver-corrected mass when such a correction is
available and otherwise the raw log-quadrature mass.  This mass normalization is
what stabilizes the reconstructed peak height after the location and width have
been determined.

For the dual trait-grid implementation, let \(\Theta_u\) be the fine phase grid
and \(\Theta_W\) the coarse amplitude grid.  The reconstruction first computes
\[
    W_U=I_{W\to u}W_c,
\]
where \(I_{W\to u}\) is a periodic interpolation operator, and then applies the
log-stabilized quadrature on \(\Theta_u\):
\[
    \rho_j=\sum_{\theta_k\in\Theta_u}\omega_k
    \exp\left(\log\bigl(\max\{(W_U)_{j,k},W_{\rm floor}\}\bigr)
    +\frac{u_k}{\varepsilon}\right).
\]
The phase equation and the peak localization are therefore resolved on the fine
one-dimensional grid, whereas the space-dependent amplitude is stored and
advanced on the coarse grid.  Since \(W\) carries all spatial degrees of
freedom, this reduction is especially useful when the spatial discretization is
large or high-dimensional.


\section{Density-compatible variant and a sufficient condition for the one-sided remainder bound}
\label{app:density_compatible}

We record a density-compatible variant of the semi-discrete WKB formulation.
This variant is not needed in the statement of the main scheme, but it gives a
concrete setting in which the weighted remainder condition in
Assumption~\ref{ass:remainder} can be verified by discrete summation by parts.

The phase equation is kept unchanged.  The amplitude equation is replaced by
\begin{equation}
\label{eq:density_compatible_W}
\begin{aligned}
&\varepsilon
\frac{\mathrm d}{\mathrm dt}W_{j,k}^\varepsilon
-
D_k\delta_x^2W_{j,k}^\varepsilon
-
\varepsilon^2(E_k^\varepsilon)^{-1}
\delta_\theta^2
\bigl(W_{j,\cdot}^\varepsilon E_\cdot^\varepsilon\bigr)_k
\\
&\qquad =
W_{j,k}^\varepsilon
\bigl(
K_j-\rho_{j,Q}^\varepsilon
+\mathcal H_k(\boldsymbol\rho_Q)
+\widehat H(D_\theta^-u_k^\varepsilon,D_\theta^+u_k^\varepsilon)
-\varepsilon\delta_\theta^2u_k^\varepsilon
\bigr),
\end{aligned}
\end{equation}
where
\[
\delta_\theta^2
\bigl(W_{j,\cdot}^\varepsilon E_\cdot^\varepsilon\bigr)_k
=
\frac{
W_{j,k+1}^\varepsilon E_{k+1}^\varepsilon
-2W_{j,k}^\varepsilon E_k^\varepsilon
+W_{j,k-1}^\varepsilon E_{k-1}^\varepsilon
}{(\Delta\theta)^2}.
\]
Multiplying \eqref{eq:density_compatible_W} by \(E_k^\varepsilon\) and using
the phase equation gives
\begin{equation}
\label{eq:density_compatible_n}
\varepsilon\frac{\mathrm d}{\mathrm dt} n_{j,k}^\varepsilon
-
D_k\delta_x^2 n_{j,k}^\varepsilon
-
\varepsilon^2\delta_\theta^2 n_{j,k}^\varepsilon
=
\bigl(K_j-\rho_{j,Q}^\varepsilon\bigr)n_{j,k}^\varepsilon .
\end{equation}
Thus the reconstructed density satisfies a conservative semi-discrete analogue
of the original density equation.  In the notation of
\eqref{eq:semi_discrete_n}, the remainder is exactly
\[
    R_{j,k}^\varepsilon
    =
    \varepsilon^2\delta_\theta^2n_{j,k}^\varepsilon .
\]

Consequently, the weighted remainder satisfies
\[
    R_{j}^{\varepsilon,D}
    =
    \varepsilon^2\Delta\theta
    \sum_k
    \frac{\delta_\theta^2n_{j,k}^\varepsilon}{D_k}.
\]
Using periodic summation by parts in the trait variable, we obtain
\[
    R_{j}^{\varepsilon,D}
    =
    \varepsilon^2\Delta\theta
    \sum_k
    n_{j,k}^\varepsilon
    \delta_\theta^2\left(\frac{1}{D_k}\right).
\]
Hence, for \(0<\varepsilon\le 1\),
\[
    \bigl(R_{j}^{\varepsilon,D}\bigr)_+
    \le
    C_{D,h} m_j^\varepsilon,
    \qquad
    C_{D,h}
    =
    \max_k
    \left|
    \delta_\theta^2\left(\frac{1}{D_k}\right)
    \right|,
\]
where
\[
    m_j^\varepsilon
    =
    \Delta\theta\sum_k n_{j,k}^\varepsilon .
\]
This verifies Assumption~\ref{ass:remainder} with
\(C_R=C_{D,h}\) and \(r_h=0\).  If \(D\) is smooth and a standard centered
periodic difference is used, \(C_{D,h}\) remains bounded under regular
trait-grid refinement.



\section{Direct density solver and two-dimensional finite-element realization}\label{app:direct_2d}

For comparison, the original density equation can be discretized by the old linear implicit update
\begin{equation}\label{eq:direct_density_fully_new}
    \varepsilon\frac{n_{j,k}^{n+1}-n_{j,k}^n}{\tau}
    -D_k\delta_x^2n_{j,k}^{n+1}
    -\varepsilon^2\delta_\theta^2n_{j,k}^{n+1}
    =n_{j,k}^{n+1}(K_j-\rho_j^n),
    \qquad
    \rho_j^n=\Delta\theta\sum_k n_{j,k}^n.
\end{equation}
{\color{red} In the revised code, the direct density solver also accepts the Patankar reaction option
\begin{equation}\label{eq:direct_density_patankar_new}
    \varepsilon\frac{n_{j,k}^{n+1}-n_{j,k}^n}{\tau}
    -D_k\delta_x^2n_{j,k}^{n+1}
    -\varepsilon^2\delta_\theta^2n_{j,k}^{n+1}
    =(K_j-\rho_j^n)_+n_{j,k}^n-(K_j-\rho_j^n)_-n_{j,k}^{n+1}.
\end{equation}
This option is used to match the positivity-oriented WKB update when direct and WKB runs are compared.} The same spatial and trait grids are used for the direct and WKB solvers in the one-dimensional comparison tests.

For the two-dimensional spatial experiment, the spatial Laplacian and the eigenvalue problem are discretized by linear finite elements on a triangular mesh.  For each trait node \(\theta_k\), the discrete effective Hamiltonian is computed from the generalized eigenvalue problem
\begin{equation}\label{eq:fem_H_new}
    A_k(\rho_Q)N^{(k)}=H_k(\rho_Q)M N^{(k)},
\end{equation}
where \(M\) is the finite-element mass matrix and
\begin{equation}\label{eq:fem_A_new}
    (A_k)_{ab}
    =D_k\int_\Omega \nabla\varphi_a\cdot\nabla\varphi_b\,dx
    -\int_\Omega (K(x)-\rho_Q(x))\varphi_a\varphi_b\,dx.
\end{equation}
The Neumann boundary condition is natural in this finite-element formulation.  The phase equation remains one-dimensional in \(\theta\), and the amplitude equation is solved for all spatial finite-element degrees of freedom and trait nodes.
\end{document}